\documentclass[10pt,a4paper]{amsart}
\usepackage{amsmath,amssymb,amsthm,graphicx}
\usepackage{verbatim}
\newenvironment{reference}{\comment}{\endcomment}
%\newenvironment{reference}{}{}
\newenvironment{slogan}{\comment}{\endcomment}
\newenvironment{history}{\comment}{\endcomment}

% For commutative diagrams we use Xy-pic
\usepackage[all]{xy}

% We use 2cell for 2-commutative diagrams.
\xyoption{2cell}
% Theorem environments.
%
\theoremstyle{plain}
\newtheorem{theorem}[subsection]{Theorem}
\newtheorem{proposition}[subsection]{Proposition}
\newtheorem{lemma}[subsection]{Lemma}

\theoremstyle{definition}
\newtheorem{definition}[subsection]{Definition}
\newtheorem{example}[subsection]{Example}
\newtheorem{exercise}[subsection]{Exercise}
\newtheorem{situation}[subsection]{Situation}

\theoremstyle{remark}
\newtheorem{remark}[subsection]{Remark}
\newtheorem{remarks}[subsection]{Remarks}

\numberwithin{equation}{subsection}

% Macros
%
\def\lim{\mathop{\mathrm{lim}}\nolimits}
\def\colim{\mathop{\mathrm{colim}}\nolimits}
\def\Spec{\mathop{\mathrm{Spec}}}
\def\Hom{\mathop{\mathrm{Hom}}\nolimits}
\def\Ext{\mathop{\mathrm{Ext}}\nolimits}
\def\SheafHom{\mathop{\mathcal{H}\!\mathit{om}}\nolimits}
\def\SheafExt{\mathop{\mathcal{E}\!\mathit{xt}}\nolimits}
\def\Sch{\mathit{Sch}}
\def\Mor{\mathop{\mathrm{Mor}}\nolimits}
\def\Ob{\mathop{\mathrm{Ob}}\nolimits}
\def\Sh{\mathop{\mathit{Sh}}\nolimits}
\def\NL{\mathop{N\!L}\nolimits}
\def\CH{\mathop{\mathrm{CH}}\nolimits}
\def\proetale{{pro\text{-}\acute{e}tale}}
\def\etale{{\acute{e}tale}}
\def\QCoh{\mathit{QCoh}}
\def\Ker{\mathop{\mathrm{Ker}}}
\def\Im{\mathop{\mathrm{Im}}}
\def\Coker{\mathop{\mathrm{Coker}}}
\def\Coim{\mathop{\mathrm{Coim}}}

% Boxtimes
%
\DeclareMathSymbol{\boxtimes}{\mathbin}{AMSa}{"02}

%
% Macros for moduli stacks/spaces
%
\def\QCohstack{\mathcal{QC}\!\mathit{oh}}
\def\Cohstack{\mathcal{C}\!\mathit{oh}}
\def\Spacesstack{\mathcal{S}\!\mathit{paces}}
\def\Quotfunctor{\mathrm{Quot}}
\def\Hilbfunctor{\mathrm{Hilb}}
\def\Curvesstack{\mathcal{C}\!\mathit{urves}}
\def\Polarizedstack{\mathcal{P}\!\mathit{olarized}}
\def\Complexesstack{\mathcal{C}\!\mathit{omplexes}}
% \Pic is the operator that assigns to X its picard group, usage \Pic(X)
% \Picardstack_{X/B} denotes the Picard stack of X over B
% \Picardfunctor_{X/B} denotes the Picard functor of X over B
\def\Pic{\mathop{\mathrm{Pic}}\nolimits}
\def\Picardstack{\mathcal{P}\!\mathit{ic}}
\def\Picardfunctor{\mathrm{Pic}}
\def\Deformationcategory{\mathcal{D}\!\mathit{ef}}
\usepackage{cleveref}
\crefname{originalreference}{original reference}{original references}
\Crefname{originalreference}{Original reference}{Original references}
\makeatletter
\newlabel{remark-trace}{{0FIK}{source}{}{sourcepointer-1}{}}
\newlabel{remark-trace@cref}{{[originalreference][0][]0FIK}{[1][0][]source}{}{}{}}
\makeatother
\title{The degree of the unit of a Galois extension}
\author{The Stacks Project Authors}
\date{stacks-project:weil.tex:local-snapshot-20261012; GNU FDL 1.2 or later}
\begin{document}
\maketitle
\noindent\textit{Editorial note.} All delivered mathematical text is the author's own, copied from the pinned Stacks Project chapter file weil.tex; no statement or proof step was rewritten, reordered or removed.\par\medskip
\noindent\textit{Editorial note.} The counted proof is the author's complete argument: the decomposition of the diagonal into graphs of the automorphisms of the Galois extension, and the computation of the pairing counting the automorphisms.\par\medskip
\noindent\textit{Editorial note.} The label lemma-integrate-1 names the counted statement, and the proof environment immediately following it in the source is delivered verbatim as its proof.\par\medskip
\noindent\textit{Editorial note.} The counted statement is carried by the author's own lemma environment in the chapter source.\par\medskip
\noindent\textit{Editorial note.} Pre-open gate: python operations/agent-authoring/tools/precard\_gate.py --source weil.tex --label lemma-integrate-1 --start 5063 --end 5100 --local-id 28151 returned RESULT=OK (L28151 free, occupancy FREE, public interval clear). Pre-assembly tag self-check: canonical tag resolved and not present in reference/public-items.jsonl.\par\medskip
\noindent\textit{Editorial note.} The statement and proof are copied from the pinned chapter file; the cited inputs are kept as external pointers with their canonical Stacks tags.\par\medskip
\noindent\textit{Allowed standard prerequisites (editorial).}\begin{itemize}
\item The diagonal of the spectrum of a Galois extension splits into the graphs of the automorphisms. The decomposition of the product of the spectrum of a finite Galois extension with itself into the graphs of the elements of the Galois group. Writes the fundamental class of the product as a sum of the pullbacks of the class of the diagonal.
\item The cycle class of a nonempty scheme is one. Same chapter, the lemma that the cycle class of the fundamental class of a nonempty scheme is one. Shows that the class of the identity in the degree zero cohomology pulls back to the sum of the pullbacks of the class of the diagonal.
\item The degree map is the pairing determined by the diagonal. Same chapter, the remark identifying the degree map with the map from the axiom defining it through the class of the diagonal. Identifies the computed pairing with the degree of the unit.
\end{itemize}\medskip
\ifdefined\hypertarget\hypertarget{sourcepointer-1}{}\fi
\noindent\textit{Source reference.} Original reference 0FIK: sources/stacks-weil/pinned-original-source.tex, source lines 4929--4929. Prerequisite pointer invoked at proof line(s) 5099; it names a standard fact of the same chapter that the counted argument uses, and every step of the counted proof built on it is retained verbatim in the delivered body.\par\medskip
\begin{lemma}
\label{lemma-integrate-1}
Assume given data (D0), (D1), and (D2') satisfying axioms (A1) -- (A7).
If $k'/k$ is a Galois extension, then we have
$\int_{\Spec(k')} 1 = [k' : k]$.
\end{lemma}

\begin{proof}
We observe that
$$
\Spec(k') \times \Spec(k') =
\coprod\nolimits_{\sigma \in \text{Gal}(k'/k)}
(\Spec(\sigma) \times \text{id})^{-1} \Delta
$$
as cycles on $\Spec(k') \times \Spec(k')$.
Our construction of $\gamma$ always sends $[X]$ to $1$ in $H^0(X)$. Thus
$1 \otimes 1 = 1 = \sum (\Spec(\sigma) \times \text{id})^*\gamma([\Delta])$.
Denote $\lambda : H^0(\Spec(k')) \to F$ the map from
axiom (A6), in other words $(\text{id} \otimes \lambda)(\gamma(\Delta)) = 1$
in $H^0(\Spec(k'))$. We obtain
\begin{align*}
\lambda(1) 1
& =
(\text{id} \otimes \lambda)(1 \otimes 1) \\
& =
(\text{id} \otimes \lambda)(
\sum (\Spec(\sigma) \times \text{id})^*\gamma([\Delta])) \\
& =
\sum (\Spec(\sigma) \times \text{id})^*(
(\text{id} \otimes \lambda)(\gamma([\Delta])) \\
& =
\sum (\Spec(\sigma) \times \text{id})^*(1) \\
& =
[k' : k]
\end{align*}
Since $\lambda$ is another name for $\int_{\Spec(k')}$
(Remark \ref{remark-trace}) the proof is complete.
\end{proof}

\end{document}
